\hojaejercicio{Método de Gauss Jordan}

% Ejercicio 1
\ejercicio{resuelto}{profesor}{
    Sea $A \in M_n$ una matriz invertible.
    \begin{itemize}
        \item[a)] Adaptar el método de Gauss Jordan para calcular $\widetilde{M}$ que cumpla que $\widetilde{M}A=I$.
        \item[b)] Usando el apartado anterior, explicar como calcular $A^{-1}$ con el método de Gauss Jordan.
    \end{itemize}
}{
    Veamos ambas partes por separado:
    \begin{itemize}
        \item[a)] Primero, según el método de Gauss Jordan, podemos realizar operaciones elementales sobre las filas de $A$ para llevarla a la forma escalonada reducida, es decir, encontrar matrices $E$ y $T$ (triangular superior) tal que $EA=T$. Luego, podemos aplicar operaciones elementales sobre las filas de $T$ para llevarla a la forma identidad, es decir, encontrar matrices $R$ y $I$ tal que $RT=I$. Combinando ambas partes, tenemos que $R(EA)=I$, lo que implica que $\widetilde{M}=RE$ cumple la condición $\widetilde{M}A=I$.\\
        Para el cálculo de $R$, podremos ir descompiniendola como el producto de $n$ matrices elementales, cada una de las cuales corresponde a una operación elemental sobre las filas de $T$, y de manera análoga al método de Gauss, vamos a ir aplicando estas operaciones sobre la última columna de $T$, dejando una submatriz con la que repetir el proceso:
        \[
            R_1 \underbrace{T_0}_{=T} =
            \left(
            \begin{array}{c|c}
                I & - \beta \\
                \hline
                0 & 1
            \end{array}
            \right)
            \left(
            \begin{array}{c|c}
                T_1 & \beta \\
                \hline
                0 & 1
            \end{array}
            \right) =
            \left(
            \begin{array}{c|c}
                T_1 & 0 \\
                \hline
                0 & 1
            \end{array}
            \right)
        \]
        Y podemos repetir el proceso con $T_1$ y así sucesivamente hasta llegar a la forma identidad, obteniendo $R$ como el producto de las matrices elementales correspondientes a cada operación elemental sobre las filas de $T$.\\
        \item[b)] Para calcular $A^{-1}$ con el método de Gauss Jordan, podemos aplicar el método de Gauss Jordan a la matriz, primero escogiendo un numero distinto de cero en cada columna, de la forma $\alpha_{ij}$ donde $ i \geq j$, que debe existir al ser $A$ invertible, y luego eliminando los elementos de esa columna, y asi sucesivamente hasta llegar a la forma identidad. Veamos una iteración del proceso:
        \[
            \left(
            \begin{array}{c|c|c}
                I & \alpha_1 & \beta_1 \\
                \hline
                0 & 1 & \beta_2 \\
                \hline
                0 & \alpha_2 & \beta_3
            \end{array}
            \right)
            \xrightarrow{R_i}
            \left(
            \begin{array}{c|c|c}
                I & 0 & \beta'_1 \\
                \hline
                0 & 1 & \beta'_2 \\
                \hline
                0 & 0 & \beta'_3
            \end{array}
            \right)
        \]
        Notese que si realizamos este proceso sobre la matriz aumentada $[A | I]$, al llegar a la forma identidad en la parte izquierda, la parte derecha de la matriz aumentada se habrá transformado en $A^{-1}$, es decir, $[I | A^{-1}]$.\\
        Esto es debido a que las operaciones elementales sobre las filas de $A$ se corresponden con multiplicar $A$ por matrices elementales, y estas mismas operaciones aplicadas a la parte derecha de la matriz aumentada se corresponden con multiplicar $I$ por las mismas matrices elementales, lo que da como resultado $A^{-1}$ al llegar a la forma identidad en la parte izquierda, ya que al haber multiplicado $A$ por las matrices elementales correspondientes a cada operación elemental sobre las filas de $A$, se habrá obtenido la forma identidad, lo que implica que $I$ se habrá multiplicado por las mismas matrices elementales, dando como resultado $A^{-1} = R_n \cdots R_1 E A = R_n \cdots R_1 E = \widetilde{M}$.\\
        En resumen, para calcular $A^{-1}$ con el método de Gauss Jordan, podemos aplicar el método de Gauss Jordan a la matriz aumentada $[A | I]$, realizando operaciones elementales sobre las filas de $A$ para llevarla a la forma identidad, y aplicando las mismas operaciones elementales sobre las filas de $I$ para obtener $A^{-1}$ en la parte derecha de la matriz aumentada.
        \end{itemize}     
}

% Ejercicio 2
\ejercicio{enunciado}{no}{}{}

% Ejercicio 3
\ejercicio{enunciado}{no}{}{}

% Ejercicio 4
\ejercicio{enunciado}{no}{}{}

% Ejercicio 5
\ejercicio{enunciado}{no}{}{}

% Ejercicio 6
\ejercicio{enunciado}{no}{}{}

% Ejercicio 7
\ejercicio{resuelto}{profesor}{
    Algoritmo para la resolución de sistemas tridiagonales. Se considera el sistema lineal $Ax = d$ para una matriz tridiagonal
    \[
        A = \begin{pmatrix}
            b_1 & c_1 & & & \\
            a_2 & b_2 & c_2 & & \\
            & \ddots & \ddots & \ddots & \\
            & & a_{n-1} & b_{n-1} & c_{n-1} \\
            & & & a_n & b_n
        \end{pmatrix}.
    \]
    \begin{itemize}
        \item[a)] Llamando $\delta_0 = 1$ y $\delta_k$ al menor principal de orden $k$ de $A$ ($k = 1, 2, \ldots, n$), probar, desarrollando el determinante por la última columna, que $\delta_1 = b_1$ y
        \[
            \delta_k = b_k \delta_{k-1} - a_k c_{k-1} \delta_{k-2}, \quad k = 2, 3, \ldots, n.
        \]
        \item[b)] Definimos las sucesiones
        \[
            \begin{cases}
                m_1 = b_1 \\
                m_k = b_k - \dfrac{c_{k-1}}{m_{k-1}} a_k, \quad k = 2, 3, \ldots, n
            \end{cases}
            \qquad \text{y} \qquad
            \begin{cases}
                g_1 = \dfrac{d_1}{m_1} \\[6pt]
                g_k = \dfrac{d_k - g_{k-1} a_k}{m_k}, \quad k = 2, 3, \ldots, n.
            \end{cases}
        \]
        Probar que $m_k = \dfrac{\delta_k}{\delta_{k-1}}$ y deducir que si $\delta_k \neq 0$, $k = 1, 2, \ldots, n$, entonces $m_k \neq 0$ y, por tanto, los números $m_k$ y $g_k$, $k = 1, 2, \ldots, n$, están bien definidos.
        \item[c)] Demostrar que la solución del sistema viene dada por
        \[
            \begin{cases}
                x_n = g_n \\
                x_k = g_k - \dfrac{c_k}{m_k} x_{k+1}, \quad k = n-1, n-2, \ldots, 1.
            \end{cases}
        \]
        \item[d)] Comprobar que
        \[
            \begin{pmatrix}
                1 & & & & \\
                \dfrac{a_2}{m_1} & 1 & & & \\
                & \ddots & \ddots & & \\
                & & \dfrac{a_{n-1}}{m_{n-2}} & 1 & \\
                & & & \dfrac{a_n}{m_{n-1}} & 1
            \end{pmatrix}
            \begin{pmatrix}
                m_1 & c_1 & & & \\
                & m_2 & c_2 & & \\
                & & \ddots & \ddots & \\
                & & & m_{n-1} & c_{n-1} \\
                & & & & m_n
            \end{pmatrix}
        \]
        es la factorización $LU$ de la matriz $A$.
    \end{itemize}
}{
    Podemos aplicar la descomposición LU a la matriz tridiagonal $A$, obteniendo matrices $L$ y $U$ de la forma:
        \[
            A = LU = 
            \begin{pmatrix}
                1 & 0 & 0 & \cdots & 0 \\
                \dfrac{a_2}{m_1} & 1 & 0 & \cdots & 0 \\
                0 & \dfrac{a_3}{m_2} & 1 & \cdots & 0 \\
                \vdots & & & \ddots & \\
                0 & 0 & 0 & \dfrac{a_n}{m_{n-1}} & 1
            \end{pmatrix}
            \begin{pmatrix}
                m_1 & c_1 & 0 & \cdots & 0 \\
                0 & m_2 & c_2 & \cdots & 0 \\
                0 & 0 & m_3 & \cdots & 0 \\
                \vdots & & & \ddots & c_{n-1} \\
                0 & 0 & 0 & 0 & m_n
            \end{pmatrix}
        \]
    Donde $m_k$ se define como dicta el en el enunciado del ejercicio. Este resultado se prueba trivialmente por inducción sobre $n$, comprobando que la multiplicación de las matrices $L$ y $U$ da como resultado la matriz $A$.\\
    Ahora lo que podemos hacer es aplicar la descomposición LU al sistema lineal $Ax = d$, obteniendo:
    \[
    Ax = d \implies
    LUx = d \underbrace{\implies}_{Ux = w}
    Lw = d
    \]
    Donde $w$ es un vector que se obtiene al resolver el sistema lineal $Lw = d$. Este sistema se resuelve fácilmente por remonte.\\
    Ahora podemos definir $x_n = \frac{w_n}{m_n}$, y luego ir definiendo $x_k$ para $k = n-1, n-2, \ldots, 1$ de la forma $x_{k-1} = \frac{w_{k-1} - c_{k-1} x_k}{m_{k-1}}$. De esta forma, obtenemos la solución del sistema lineal $Ax = d$ de la forma que dicta el enunciado del ejercicio.\\
    Por último, para probar la relación entre $m_k$ y $\delta_k$, podemos observar que $\delta_k$ se puede calcular como el producto de los elementos de la diagonal de la matriz $U$, es decir, $\delta_k = m_1 m_2 \cdots m_k$. Por lo tanto, $m_k = \frac{\delta_k}{\delta_{k-1}}$, lo que implica que si $\delta_k \neq 0$ para $k = 1, 2, \ldots, n$, entonces $m_k \neq 0$ para $k = 1, 2, \ldots, n$, y por tanto, los números $m_k$ y $g_k$ están bien definidos.\\
    En resumen, el algoritmo para la resolución de sistemas tridiagonales consiste en aplicar la descomposición LU a la matriz tridiagonal $A$, obteniendo las matrices $L$ y $U$, y luego resolver el sistema lineal $Lw = d$ por remonte, y finalmente obtener la solución del sistema lineal $Ax = d$ utilizando las fórmulas dadas en el enunciado del ejercicio.
}

% Ejercicio 8
\ejercicio{enunciado}{no}{}{}

% Ejercicio 9
\ejercicio{enunciado}{no}{}{}

% Ejercicio 10
\ejercicio{enunciado}{no}{}{}

% Ejercicio 11
\ejercicio{enunciado}{no}{}{}

% Ejercicio 12
\ejercicio{enunciado}{no}{}{}

% Ejercicio 13
\ejercicio{enunciado}{no}{}{}

% Ejercicio 14
\ejercicio{enunciado}{no}{}{}

% Ejercicio 15
\ejercicio{enunciado}{no}{}{}

% Ejercicio 16
\ejercicio{enunciado}{no}{}{}

% Ejercicio 17
\ejercicio{enunciado}{no}{}{}